\documentclass[11pt,a4paper]{article}

\usepackage[utf8]{inputenc}
\usepackage[T1]{fontenc}
\usepackage{lmodern}
\usepackage[spanish,es-tabla,es-noshorthands]{babel}
\usepackage[a4paper,margin=2.3cm]{geometry}
\usepackage{xcolor}
\usepackage{listings}
\usepackage{tcolorbox}
\usepackage{enumitem}
\usepackage{titlesec}
\usepackage{parskip}
\usepackage{hyperref}

\definecolor{azulOscuro}{HTML}{1E3C78}
\definecolor{azulMedio}{HTML}{2A5BB1}
\definecolor{grisCodigo}{HTML}{F5F6FA}
\definecolor{grisBorde}{HTML}{D0D4DD}
\definecolor{kw}{HTML}{0033B3}
\definecolor{str}{HTML}{067D17}
\definecolor{com}{HTML}{8C8C8C}
\definecolor{tipo}{HTML}{8B4513}
\definecolor{verdeTAC}{HTML}{2D7D46}
\definecolor{rojoOpt}{HTML}{B23B3B}

\hypersetup{
  colorlinks=true,
  linkcolor=azulMedio,
  urlcolor=azulMedio,
  pdftitle={Compiladores 2026 - Fase 2}
}

\titleformat{\section}
  {\color{azulOscuro}\Large\bfseries}
  {\thesection.}{0.6em}{}
\titleformat{\subsection}
  {\color{azulMedio}\large\bfseries}
  {\thesubsection.}{0.5em}{}
\titleformat{\subsubsection}
  {\color{azulMedio}\normalsize\bfseries}
  {\thesubsubsection.}{0.5em}{}

\lstdefinestyle{py}{
  language=Python,
  backgroundcolor=\color{grisCodigo},
  basicstyle=\ttfamily\footnotesize,
  keywordstyle=\color{kw}\bfseries,
  stringstyle=\color{str},
  commentstyle=\color{com}\itshape,
  frame=single,
  framerule=0.4pt,
  rulecolor=\color{grisBorde},
  numbers=none,
  showstringspaces=false,
  breaklines=true,
  xleftmargin=8pt,
  xrightmargin=8pt,
  aboveskip=6pt,
  belowskip=6pt,
}
\lstdefinestyle{js}{
  language=Java,
  backgroundcolor=\color{grisCodigo},
  basicstyle=\ttfamily\footnotesize,
  keywordstyle=\color{kw}\bfseries,
  stringstyle=\color{str},
  commentstyle=\color{com}\itshape,
  frame=single,
  framerule=0.4pt,
  rulecolor=\color{grisBorde},
  numbers=none,
  showstringspaces=false,
  breaklines=true,
  xleftmargin=8pt,
  xrightmargin=8pt,
  aboveskip=6pt,
  belowskip=6pt,
}
\lstdefinestyle{tac}{
  backgroundcolor=\color{grisCodigo},
  basicstyle=\ttfamily\footnotesize,
  frame=single,
  framerule=0.4pt,
  rulecolor=\color{grisBorde},
  numbers=none,
  showstringspaces=false,
  breaklines=true,
  xleftmargin=8pt,
  xrightmargin=8pt,
  aboveskip=6pt,
  belowskip=6pt,
}

\newtcolorbox{caja}[1]{
  colback=grisCodigo,
  colframe=azulMedio,
  title=\textbf{#1},
  fonttitle=\color{white},
  boxrule=0.5pt,
  arc=2pt,
  left=6pt,right=6pt,top=4pt,bottom=4pt,
  before skip=8pt, after skip=8pt,
}

\title{\vspace{-1.2cm}\color{azulOscuro}\textbf{Compiladores 2026 --- Fase 2}\\[0.3em]
\large\color{black}Frontend/Backend, Código Intermedio, Optimización y Visualización}
\author{}
\date{}

\begin{document}
\maketitle
\vspace{-1.4cm}

\section{Resumen}

En esta fase se reorganizó el compilador en \textbf{frontend} (análisis) y
\textbf{backend} (transformación), se incorporó la generación de
\textbf{código intermedio} en formato \textit{Three-Address Code} (TAC) y
se construyó un \textbf{optimizador} que trabaja sobre ese TAC. Además se
extendió la interfaz web para mostrar el TAC original, el optimizado y la
traza del optimizador.

A continuación se presenta el código real del proyecto, archivo por archivo,
explicando qué hace cada bloque.

\section{Separación Frontend / Backend}

\subsection{Estructura física}

\begin{lstlisting}[style=tac]
files (1)/
|-- bridge.py                       # orquestador (CLI)
|-- intermedio.py                   # contrato compartido (Quad + formatter)
|-- frontend/
|   |-- __init__.py
|   |-- lexer.py
|   |-- tabla_simbolos.py
|   |-- parser.py
|   `-- generador_intermedio.py
`-- backend/
    |-- __init__.py
    `-- optimizador.py
\end{lstlisting}

\subsection{\texttt{frontend/\_\_init\_\_.py} --- qué exporta el frontend}

\begin{lstlisting}[style=py]
from .lexer              import Lexer
from .tabla_simbolos     import TablaSimbolos, Simbolo
from .parser             import build_tree, get_kids, node_label, TIPOS_DATO, RESERVED
from .generador_intermedio import GeneradorTAC

__all__ = [
    "Lexer",
    "TablaSimbolos", "Simbolo",
    "build_tree", "get_kids", "node_label", "TIPOS_DATO", "RESERVED",
    "GeneradorTAC",
]
\end{lstlisting}

\textbf{Qué hace:} agrupa todos los componentes del análisis en un solo paquete
y los re-exporta. El \texttt{bridge.py} hace
\texttt{from frontend import Lexer, build\_tree, GeneradorTAC} y no se entera
de cómo está organizado por dentro.

\subsection{\texttt{backend/\_\_init\_\_.py} --- qué exporta el backend}

\begin{lstlisting}[style=py]
from .optimizador import optimizar

__all__ = ["optimizar"]
\end{lstlisting}

\textbf{Qué hace:} expone únicamente la función \texttt{optimizar(quads)}.
El backend hoy sólo optimiza, pero la interfaz pública queda lista para
agregar generadores de código objeto en fases siguientes sin tocar el resto.

\section{Frontera compartida: \texttt{intermedio.py}}

Este archivo es el \textbf{único} módulo que ambos lados conocen. Define la
estructura de un cuádruplo y un \textit{pretty-printer} para mostrarlo.

\subsection{La dataclass \texttt{Quad}}

\begin{lstlisting}[style=py]
@dataclass
class Quad:
    op:   str
    arg1: Optional[str] = None
    arg2: Optional[str] = None
    dest: Optional[str] = None

    def as_tuple(self):
        return (self.op,
                self.arg1 if self.arg1 is not None else "_",
                self.arg2 if self.arg2 is not None else "_",
                self.dest if self.dest is not None else "_")
\end{lstlisting}

\textbf{Qué hace:}
\begin{itemize}[leftmargin=*,nosep]
  \item \texttt{op}: el operador (\texttt{+}, \texttt{=}, \texttt{goto}, \texttt{label}, \dots).
  \item \texttt{arg1}, \texttt{arg2}: los dos operandos fuente (cualquiera puede faltar).
  \item \texttt{dest}: el destino. Para asignaciones y operaciones es el
        nombre/temporal donde queda el resultado; para saltos es la etiqueta
        a la que se salta.
  \item \texttt{as\_tuple()}: devuelve la 4-tupla con \texttt{"\_"} cuando hay
        \texttt{None}, útil para imprimir o serializar.
\end{itemize}

\subsection{\texttt{formatear\_tac} --- pretty-printer para la UI}

\begin{lstlisting}[style=py]
def formatear_tac(quads):
    filas = []
    n = 0
    for q in quads:
        op, a1, a2, d = q.op, q.arg1, q.arg2, q.dest
        if op == "label":
            instr = f"{d}:"
            etiqueta = d
        elif op == "goto":
            instr = f"goto {d}"
        elif op == "if_false":
            instr = f"ifFalse {a1} goto {d}"
        elif op == "=":
            instr = f"{d} = {a1}"
        elif op == "print":
            instr = f"print {a1}"
        elif op == "param":
            instr = f"param {a1}"
        elif op == "call":
            instr = f"{d} = call {a1}, {a2}"
        elif op == "!":
            instr = f"{d} = !{a1}"
        else:
            instr = f"{d} = {a1} {op} {a2}"   # binarias
        n += 1
        filas.append({"n": n, "etiqueta": etiqueta, "instruccion": instr,
                      "op": op, "arg1": a1 or "-", "arg2": a2 or "-",
                      "dest": d  or "-"})
    return filas
\end{lstlisting}

\textbf{Qué hace:} recorre la lista de cuádruplos y produce una lista de
diccionarios listos para que el JS de la UI los pinte como tabla o como
texto plano. Cada \texttt{op} tiene su propio formato legible
(\texttt{ifFalse \$t2 goto \$L1}, \texttt{c = \$t1}, \dots) sin perder
los campos crudos (\texttt{op}, \texttt{arg1}, \texttt{arg2}, \texttt{dest}).

\section{Generación del TAC: \texttt{frontend/generador\_intermedio.py}}

\subsection{Mapa de operadores}

\begin{lstlisting}[style=py]
_BIN_TIPOS = {
    "MAS": "+", "MENOS": "-", "MULT": "*", "DIV": "/", "MOD": "%",
    "IGUAL": "==", "DIFERENTE": "!=",
    "MENOR": "<", "MAYOR": ">",
    "MENOR_IGUAL": "<=", "MAYOR_IGUAL": ">=",
    "Y_LOGICO": "&&", "O_LOGICO": "||",
}
\end{lstlisting}

\textbf{Qué hace:} traduce el \texttt{tipo} del token (como lo emite el lexer)
al símbolo que va a quedar dentro del cuádruplo. Por ejemplo, el token
\texttt{MAS} se vuelve \texttt{"+"} en el TAC.

\subsection{Estado del generador y helpers}

\begin{lstlisting}[style=py]
class GeneradorTAC:
    def __init__(self):
        self.codigo: list[Quad] = []
        self._t = 0
        self._l = 0

    def _nuevo_temp(self) -> str:
        self._t += 1
        return f"$t{self._t}"

    def _nueva_etiq(self) -> str:
        self._l += 1
        return f"$L{self._l}"

    def emit(self, op, arg1=None, arg2=None, dest=None):
        self.codigo.append(Quad(op, arg1, arg2, dest))
\end{lstlisting}

\textbf{Qué hace:}
\begin{itemize}[leftmargin=*,nosep]
  \item \texttt{self.codigo}: lista de cuádruplos que se va construyendo.
  \item \texttt{\_nuevo\_temp()}: devuelve nombres frescos
        \texttt{\$t1, \$t2, \dots}\ para guardar resultados intermedios.
  \item \texttt{\_nueva\_etiq()}: devuelve nombres frescos
        \texttt{\$L1, \$L2, \dots}\ para destinos de salto.
  \item \texttt{emit()}: empuja un cuádruplo al final de la lista.
\end{itemize}

\subsection{Punto de entrada y \textit{dispatcher}}

\begin{lstlisting}[style=py]
def generar(self, ast) -> list[Quad]:
    self.codigo = []
    self._t = 0
    self._l = 0
    if ast:
        self._gen(ast)
    return self.codigo

def _gen(self, node):
    if not node:
        return None
    t = node.get("type", "")
    method = getattr(self, f"_g_{t}", self._g_default)
    return method(node)
\end{lstlisting}

\textbf{Qué hace:} \texttt{generar} resetea el estado y arranca el recorrido.
\texttt{\_gen} es el \textbf{dispatcher}: mira el campo \texttt{type} del
nodo del AST y llama al método \texttt{\_g\_<Tipo>}, p.\,ej.\
\texttt{\_g\_If} para un nodo \texttt{If}. Es un visitor implementado con
\texttt{getattr}, sin \texttt{if/elif} repetitivos.

\subsection{Expresiones binarias --- patrón post-orden}

\begin{lstlisting}[style=py]
if t == "BinaryOp":
    izq = self._g_expr(node["left"])    # subarbol izquierdo
    der = self._g_expr(node["right"])   # subarbol derecho
    tmp = self._nuevo_temp()
    op_tipo = node["op"]["tipo"]
    op = _BIN_TIPOS.get(op_tipo, node["op"]["valor"])
    self.emit(op, izq, der, tmp)
    return tmp
\end{lstlisting}

\textbf{Qué hace:} traduce \texttt{a op b} en TAC. Primero genera el código
del operando izquierdo y obtiene el lugar (variable o temporal) donde
quedó su resultado; lo mismo con el derecho; pide un temporal nuevo y
emite el cuádruplo \texttt{tmp = izq op der}. El temporal devuelto se
encadena en operaciones más grandes: \texttt{(a+b)*c} produce
\texttt{\$t1 = a + b}\ y luego \texttt{\$t2 = \$t1 * c}.

\subsection{Estructuras de control}

\subsubsection{If / else}

\begin{lstlisting}[style=py]
def _g_If(self, node):
    cond = self._g_expr(node.get("cond"))
    L_else = self._nueva_etiq()
    L_end  = self._nueva_etiq() if node.get("elseBody") else L_else

    self.emit("if_false", cond, None, L_else)
    self._gen(node.get("body"))
    if node.get("elseBody"):
        self.emit("goto", None, None, L_end)
        self.emit("label", None, None, L_else)
        self._gen(node["elseBody"])
        self.emit("label", None, None, L_end)
    else:
        self.emit("label", None, None, L_else)
\end{lstlisting}

\textbf{Qué hace:} si la condición es falsa salta a \texttt{L\_else};
si no, ejecuta el cuerpo. Cuando hay \texttt{else}, después del cuerpo se
hace \texttt{goto L\_end} para no caer en el bloque \texttt{else}.
Las dos etiquetas marcan los puntos de aterrizaje.

\subsubsection{While}

\begin{lstlisting}[style=py]
def _g_While(self, node):
    L_inicio = self._nueva_etiq()
    L_fin    = self._nueva_etiq()
    self.emit("label", None, None, L_inicio)
    cond = self._g_expr(node.get("cond"))
    self.emit("if_false", cond, None, L_fin)
    self._gen(node.get("body"))
    self.emit("goto", None, None, L_inicio)
    self.emit("label", None, None, L_fin)
\end{lstlisting}

\textbf{Qué hace:} pone una etiqueta al inicio del lazo, evalúa la condición,
salta al final si es falsa, ejecuta el cuerpo y vuelve al inicio con un
\texttt{goto}. \texttt{L\_fin} es el destino tanto del \texttt{ifFalse}
como del eventual \textit{break} (no implementado en esta fase).

\subsubsection{Print}

\begin{lstlisting}[style=py]
def _g_Print(self, node):
    arg = self._g_expr(node.get("arg"))
    self.emit("print", arg, None, None)
\end{lstlisting}

\textbf{Qué hace:} genera el código para evaluar el argumento (puede ser
una expresión compleja) y emite un cuádruplo \texttt{print arg}.

\subsection{Ejemplo de salida}

Código fuente:

\begin{lstlisting}[style=tac]
int a = 3;
int b = 4;
int c = a + b;
if (c > 5) { print(c); } else { print(0); }
\end{lstlisting}

TAC producido por \texttt{GeneradorTAC}:

\begin{lstlisting}[style=tac,basicstyle=\ttfamily\footnotesize\color{verdeTAC}]
1:  a = 3
2:  b = 4
3:  $t1 = a + b
4:  c = $t1
5:  $t2 = c > 5
6:  ifFalse $t2 goto $L1
7:  print c
8:  goto $L2
9:  $L1:
10: print 0
11: $L2:
\end{lstlisting}

\section{Optimización: \texttt{backend/optimizador.py}}

\subsection{Tabla de operaciones evaluables en compile-time}

\begin{lstlisting}[style=py]
_BIN_OPS = {
    "+":  lambda a, b: a + b,
    "-":  lambda a, b: a - b,
    "*":  lambda a, b: a * b,
    "/":  lambda a, b: a / b if isinstance(a, float) or isinstance(b, float) else a // b,
    "%":  lambda a, b: a % b,
    "==": lambda a, b: a == b,
    "!=": lambda a, b: a != b,
    "<":  lambda a, b: a < b,   ">":  lambda a, b: a > b,
    "<=": lambda a, b: a <= b,  ">=": lambda a, b: a >= b,
    "&&": lambda a, b: bool(a) and bool(b),
    "||": lambda a, b: bool(a) or bool(b),
}
\end{lstlisting}

\textbf{Qué hace:} mapea cada operador del TAC a su implementación Python.
Lo usa el \textit{constant folding} para evaluar la operación si los dos
argumentos resultan ser constantes literales.

\subsection{Pasada 1 --- Constant Folding y simplificación algebraica}

\begin{lstlisting}[style=py]
def _constant_folding(codigo):
    nuevo = []
    cambio = False
    for q in codigo:
        if q.op in _BIN_OPS:
            a1, a2 = q.arg1, q.arg2
            # Identidades algebraicas
            if q.op == "+" and a1 == "0":
                nuevo.append(Quad("=", a2, None, q.dest)); cambio = True; continue
            if q.op == "+" and a2 == "0":
                nuevo.append(Quad("=", a1, None, q.dest)); cambio = True; continue
            if q.op == "-" and a2 == "0":
                nuevo.append(Quad("=", a1, None, q.dest)); cambio = True; continue
            if q.op == "*" and (a1 == "1"):
                nuevo.append(Quad("=", a2, None, q.dest)); cambio = True; continue
            if q.op == "*" and (a2 == "1"):
                nuevo.append(Quad("=", a1, None, q.dest)); cambio = True; continue
            if q.op == "*" and (a1 == "0" or a2 == "0"):
                nuevo.append(Quad("=", "0", None, q.dest)); cambio = True; continue
            # Folding numerico
            if _es_numero(a1) and _es_numero(a2):
                r = _BIN_OPS[q.op](_to_num(a1), _to_num(a2))
                nuevo.append(Quad("=", _fmt(r), None, q.dest))
                cambio = True; continue
        nuevo.append(q)
    return nuevo, cambio
\end{lstlisting}

\textbf{Qué hace:} para cada cuádruplo binario, primero detecta
\textit{identidades algebraicas} (\texttt{x+0}, \texttt{x*1},
\texttt{x*0}\dots) y reemplaza la operación por una asignación directa.
Después, si los dos operandos son constantes numéricas, calcula el
resultado y reemplaza el cuádruplo por \texttt{dest = <resultado>}.
Devuelve el nuevo código y un flag \texttt{cambio}.

\subsection{Pasada 2 --- Constant / Copy Propagation}

\begin{lstlisting}[style=py]
def _propagation(codigo):
    nuevo = []
    env = {}             # variable/temp -> valor conocido
    cambio = False

    def reemplazar(x):
        return env.get(x, x) if x is not None else x

    for q in codigo:
        if q.op == "label":
            env = {}                 # se invalida en bordes de bloques
            nuevo.append(q); continue
        if q.op in ("goto", "if_false", "if_true"):
            nuevo.append(Quad(q.op, reemplazar(q.arg1), reemplazar(q.arg2), q.dest))
            env = {}; continue
        if q.op == "=":
            valor = reemplazar(q.arg1)
            nuevo.append(Quad("=", valor, None, q.dest))
            env[q.dest] = valor      # ahora dest "es" valor
            continue
        # Resto: reemplazar argumentos
        a1 = reemplazar(q.arg1); a2 = reemplazar(q.arg2)
        nuevo.append(Quad(q.op, a1, a2, q.dest))
        if q.dest is not None:
            env.pop(q.dest, None)
    return nuevo, cambio
\end{lstlisting}

\textbf{Qué hace:} mantiene un diccionario \texttt{env} con los valores que
conoce de cada variable/temporal. Cuando ve \texttt{x = 5}, recuerda
\texttt{env["x"] = "5"}. Cuando ve un uso posterior de \texttt{x}, lo
reemplaza por \texttt{"5"}. El entorno se \textbf{invalida} al cruzar una
etiqueta o un salto, porque ahí pueden venir valores de otros caminos.
Cualquier reasignación de \texttt{dest} también borra entradas viejas
asociadas a ese destino.

\subsection{Pasada 3 --- Branch Pruning}

\begin{lstlisting}[style=py]
def _branch_pruning(codigo):
    nuevo = []
    cambio = False
    for q in codigo:
        if q.op == "if_false":
            if q.arg1 in ("true", "True", "1"):
                cambio = True; continue                   # ifFalse true   -> elimina
            if q.arg1 in ("false", "False", "0"):
                nuevo.append(Quad("goto", None, None, q.dest))
                cambio = True; continue                   # ifFalse false  -> goto
        if q.op == "if_true":
            if q.arg1 in ("false", "False", "0"):
                cambio = True; continue
            if q.arg1 in ("true", "True", "1"):
                nuevo.append(Quad("goto", None, None, q.dest))
                cambio = True; continue
        nuevo.append(q)
    return nuevo, cambio
\end{lstlisting}

\textbf{Qué hace:} cuando la pasada de \textit{propagation} dejó la condición
de un salto convertida en literal (\texttt{true}/\texttt{false}), aquí se
elimina el salto inútil o se convierte en un \texttt{goto} incondicional.

\subsection{Pasada 4 --- Dead-Code Elimination}

\begin{lstlisting}[style=py]
def _dead_code(codigo):
    leidos = set()
    for q in codigo:
        for a in (q.arg1, q.arg2):
            if a is not None: leidos.add(a)
        if q.op in ("goto", "if_false", "if_true") and q.dest:
            leidos.add(q.dest)

    nuevo = []
    cambio = False
    for q in codigo:
        if q.op == "=" or q.op in _BIN_OPS or q.op == "!":
            if q.dest and _es_temp(q.dest) and q.dest not in leidos:
                cambio = True; continue           # temporal no leido => fuera
        nuevo.append(q)
    return nuevo, cambio
\end{lstlisting}

\textbf{Qué hace:} primero escanea \emph{todo} el código y arma el conjunto
de nombres que aparecen como \textbf{leídos} (en \texttt{arg1} o
\texttt{arg2}, o como destino de un salto). Luego elimina los cuádruplos
que escriben en un \textbf{temporal} (\texttt{\$tn}) que nunca se leyó.
Sólo se borran temporales: las variables del usuario se conservan aunque
parezcan muertas, porque pueden tener efecto observable
(\texttt{print} posterior, scope externo, etc.).

\subsection{Pasada 5 --- Jump Threading}

\begin{lstlisting}[style=py]
def _jump_threading(codigo):
    nuevo = []
    cambio = False
    i = 0
    while i < len(codigo):
        q = codigo[i]
        if (q.op == "goto"
                and i + 1 < len(codigo)
                and codigo[i + 1].op == "label"
                and codigo[i + 1].dest == q.dest):
            cambio = True; i += 1; continue   # goto L seguido de label L: borra goto
        nuevo.append(q); i += 1

    referenciadas = set()                     # etiquetas todavia usadas
    for q in nuevo:
        if q.op in ("goto","if_false","if_true") and q.dest:
            referenciadas.add(q.dest)
    final = []
    for q in nuevo:
        if q.op == "label" and q.dest not in referenciadas:
            cambio = True; continue           # label huerfana
        final.append(q)
    return final, cambio
\end{lstlisting}

\textbf{Qué hace:} dos limpiezas. Primero, si encuentra un
\texttt{goto L} inmediatamente seguido de \texttt{L:}, borra el goto
porque el flujo cae igual a la etiqueta. Segundo, calcula qué etiquetas
quedaron \emph{referenciadas} y elimina las que ya no apunta nadie.

\subsection{Driver: \texttt{optimizar()} --- punto fijo}

\begin{lstlisting}[style=py]
def optimizar(codigo, max_iter: int = 12):
    actual = deepcopy(codigo)
    traza  = []
    pasadas = [
        ("Constant Folding & Algebraic", _constant_folding),
        ("Constant / Copy Propagation",  _propagation),
        ("Branch Pruning",               _branch_pruning),
        ("Dead-Code Elimination",        _dead_code),
        ("Jump Threading",               _jump_threading),
    ]

    for it in range(max_iter):
        cambio_global = False
        for nombre, fn in pasadas:
            antes = len(actual)
            actual, cambio = fn(actual)
            despues = len(actual)
            if cambio:
                traza.append({"iter": it+1, "pasada": nombre,
                              "antes": antes, "despues": despues,
                              "delta": despues - antes})
                cambio_global = True
        if not cambio_global:
            break
    return actual, traza
\end{lstlisting}

\textbf{Qué hace:} aplica las cinco pasadas en orden, una y otra vez, hasta
que una iteración completa no haga ningún cambio (\textbf{punto fijo}) o
se llegue a \texttt{max\_iter}. Cada vez que una pasada modifica el
código, se anota una entrada en \texttt{traza} con el nombre, el tamaño
antes/después y el delta. Devuelve \texttt{(codigo\_optimizado, traza)}.

\subsection{Antes y después}

Código fuente:

\begin{lstlisting}[style=tac]
int a = 3;  int b = 4;
int c = a + b;     int d = c * 2;
int e = d - 0;     int f = e * 1;
print(f);
\end{lstlisting}

\begin{minipage}[t]{0.48\textwidth}
\textbf{TAC original}
\begin{lstlisting}[style=tac,basicstyle=\ttfamily\footnotesize\color{verdeTAC}]
a = 3
b = 4
$t1 = a + b
c = $t1
$t2 = c * 2
d = $t2
$t3 = d - 0
e = $t3
$t4 = e * 1
f = $t4
print f
\end{lstlisting}
\end{minipage}\hfill
\begin{minipage}[t]{0.48\textwidth}
\textbf{TAC optimizado}
\begin{lstlisting}[style=tac,basicstyle=\ttfamily\footnotesize\color{rojoOpt}]
a = 3
b = 4
c = 7
d = 14
e = 14
f = 14
print 14
\end{lstlisting}
\end{minipage}

\section{Orquestador: \texttt{bridge.py}}

\begin{lstlisting}[style=py]
from frontend  import Lexer, build_tree, GeneradorTAC
from intermedio import formatear_tac
from backend   import optimizar

# 1) lexer
lexer = Lexer()
tokens_info, tabla, errores = lexer.analizar(codigo)

# 2) parser -> AST -> TAC
ast   = build_tree(tokens_info)
quads = GeneradorTAC().generar(ast)
tac_filas = formatear_tac(quads)

# 3) optimizador
quads_opt, traza = optimizar(quads)
tac_opt_filas    = formatear_tac(quads_opt)

# 4) metricas + JSON
metricas = {
    "cuad_orig": len(quads), "cuad_opt": len(quads_opt),
    "reduccion_pct": round(100.0*(len(quads)-len(quads_opt))/len(quads), 1),
    "temps_orig": _contar_temps(quads), "temps_opt": _contar_temps(quads_opt),
}
print(json.dumps({
    "tokens": tokens_info, "simbolos": simbolos_json, "errores": errores,
    "tac": tac_filas, "tac_optimizado": tac_opt_filas,
    "metricas": metricas, "traza_optimizacion": traza,
}, ensure_ascii=False))
\end{lstlisting}

\textbf{Qué hace:} es el único archivo que une el frontend con el backend.
Hace el pipeline completo \emph{lexer $\to$ parser $\to$ generador TAC
$\to$ optimizador} y empaqueta todo (tokens, símbolos, errores, TAC
original, TAC optimizado, métricas y traza) como un único objeto JSON
que escribe a \texttt{stdout}. La UI lo invoca, lee el JSON y pinta los
paneles.

\section{Parte visual: pestaña de Código Intermedio en \texttt{index.php}}

\subsection{Lectura del JSON}

\begin{lstlisting}[style=js]
const tac    = data.tac           || [];
const tacOpt = data.tac_optimizado|| [];
const metr   = data.metricas      || {};
const traza  = data.traza_optimizacion || [];

renderIntermediate(tac, tacOpt, metr, traza);
\end{lstlisting}

\textbf{Qué hace:} extrae los cuatro campos del JSON producido por
\texttt{bridge.py} y se los pasa a la función de render.

\subsection{Render: tablas, pretty, métricas y traza}

\begin{lstlisting}[style=js]
function renderIntermediate(tac, tacOpt, metr, traza){
  // Tablas (op | arg1 | arg2 | dest)
  document.getElementById('ir-body-orig').innerHTML = tac.map(_irRowHtml).join('');
  document.getElementById('ir-body-opt').innerHTML  = tacOpt.map(_irRowHtml).join('');

  // Vista pretty-printer (texto plano alineado)
  document.getElementById('ir-pre-orig').textContent = _irPretty(tac);
  document.getElementById('ir-pre-opt').textContent  = _irPretty(tacOpt);
  document.getElementById('ir-cmp-orig-n').textContent = `${tac.length} cuadruplos`;
  document.getElementById('ir-cmp-opt-n').textContent  = `${tacOpt.length} cuadruplos`;

  // Traza del optimizador
  document.getElementById('ir-traza-body').innerHTML = (traza||[]).map(t => {
    const dColor = t.delta<0 ? 'var(--green)' : (t.delta>0 ? 'var(--red)' : 'var(--t3)');
    const dStr   = t.delta>0 ? `+${t.delta}` : String(t.delta);
    return `<tr>
      <td>${t.iter}</td><td>${escH(t.pasada)}</td>
      <td>${t.antes}</td><td>${t.despues}</td>
      <td style="color:${dColor}">${dStr}</td></tr>`;
  }).join('');

  // Metricas
  const co = metr.cuad_orig ?? tac.length;
  const cp = metr.cuad_opt  ?? tacOpt.length;
  document.getElementById('ir-stats').innerHTML =
    `Cuadruplos: <b>${co}</b> -> <b>${cp}</b> &middot; ` +
    `Reduccion: <b>${metr.reduccion_pct ?? 0}%</b> &middot; ` +
    `Temporales: <b>${metr.temps_orig ?? 0}</b> -> <b>${metr.temps_opt ?? 0}</b>`;
}
\end{lstlisting}

\textbf{Qué hace cada bloque:}
\begin{itemize}[leftmargin=*,nosep]
  \item \textbf{Tablas}: rellena las tablas de TAC original y optimizado
        con una fila por cuádruplo (\texttt{op}, \texttt{arg1},
        \texttt{arg2}, \texttt{dest}).
  \item \textbf{Pretty}: escribe la versión legible
        (estilo \texttt{c = \$t1} o \texttt{print x}) producida
        en Python por \texttt{formatear\_tac}.
  \item \textbf{Traza}: una fila por pasada que hizo cambios. El delta se
        pinta en verde si redujo y en rojo si creció (no sucede en las
        pasadas implementadas pero deja la columna preparada).
  \item \textbf{Métricas}: tarjeta con cuádruplos antes/después,
        porcentaje de reducción y número de temporales.
\end{itemize}

\subsection{Sub-vistas}

\begin{lstlisting}[style=js]
function switchIRSub(sub){
  document.querySelectorAll('.ir-view').forEach(v=>v.style.display='none');
  const map = {orig:'ir-view-orig', opt:'ir-view-opt',
               cmp:'ir-view-cmp',   info:'ir-view-info'};
  const target = document.getElementById(map[sub] || map.orig);
  if(target) target.style.display = (sub==='cmp') ? 'flex' : 'block';
}
\end{lstlisting}

\textbf{Qué hace:} permite alternar entre cuatro vistas dentro de la
pestaña: \texttt{Original}, \texttt{Optimizado}, \texttt{Comparación}
(lado a lado) e \texttt{Info}. Oculta todas las vistas y muestra sólo
la elegida.

\section{Conclusión}

La fase 2 deja el compilador con tres piezas bien diferenciadas:

\begin{itemize}[leftmargin=*]
  \item Un \textbf{frontend} (\texttt{lexer + parser + generador\_intermedio})
        que produce un AST y de ahí el TAC.
  \item Un \textbf{contrato} (\texttt{intermedio.py}, la dataclass
        \texttt{Quad}) que ambos lados conocen y nada más.
  \item Un \textbf{backend} (\texttt{optimizador.py}) con cinco pasadas
        clásicas que iteran hasta punto fijo.
\end{itemize}

La interfaz visual hace explícito el efecto del optimizador (tabla
original vs.\ optimizado, métricas y traza), lo que vuelve al
compilador útil también como herramienta de aprendizaje:
no sólo \emph{qué} se redujo, sino también \emph{por qué} pasada y en
qué iteración.

\end{document}
